\documentclass[preprint,12pt,a4paper]{elsarticle}
\usepackage{amssymb}
\usepackage{hyperref}
\setlength{\parindent}{0pt}

\journal{SoftwareX}

\begin{document}

\begin{frontmatter}

\title{OpenTelemetry-to-EEOAP Adapter: Transforming Agent Telemetry into Portable Operation Evidence}

\author{Bin Zhang}
\address{TODO: confirm public affiliation before external submission}
\ead{TODO: confirm support/contact email before external submission}

\begin{abstract}
Runtime telemetry helps developers observe agent execution, but trace records do
not automatically become portable operation evidence. The OpenTelemetry-to-EEOAP
Adapter is a small research-software package that converts controlled
OpenTelemetry-style agent trace fixtures into Execution Evidence and Operation
Accountability Profile compatible operation accountability statements. The
adapter selects the accountable agent span, extracts operation metadata,
resolves tool spans, checks parent-child relations, preserves provenance, emits
an EEOAP-compatible statement, and routes the output to the existing validator.
The evaluation includes two valid trace contexts and four invalid diagnostic
contexts. Both valid generated statements pass the existing validator with
ok=true and issue_count=0. Public release identifiers, DOI, GitHub Release, and
pushed tags remain TODO before formal submission.
\end{abstract}

\begin{keyword}
OpenTelemetry \sep EEOAP \sep agent telemetry \sep operation accountability \sep execution evidence \sep validation \sep reproducibility
\end{keyword}

\end{frontmatter}

\section*{Required Metadata}

\section*{Current code version}

\begin{table}[!h]
\begin{tabular}{|l|p{6.5cm}|p{6.5cm}|}
\hline
\textbf{Nr.} & \textbf{Code metadata description} & \textbf{Please fill in this column} \\
\hline
C1 & Current code version & TODO: assign release version before external submission. Draft label: otel-eeoap-adapter-v0.1.0-rc. \\
\hline
C2 & Permanent GitHub link to code/repository used for this code version & TODO: create public GitHub release URL, archive, DOI, or immutable public commit reference before formal submission. \\
\hline
C3 & Legal Code License & Apache-2.0 via root LICENSE. TODO: confirm final SoftwareX package license file expectations. \\
\hline
C4 & Code versioning system used & git \\
\hline
C5 & Software code languages, tools, and services used & Python; JSON fixtures; local agent-evidence CLI validator path; pytest. \\
\hline
C6 & Compilation requirements, operating environments \& dependencies & Python 3.11+; click; jsonschema; pydantic; pytest for scoped tests. TODO: confirm final operating-system statement. \\
\hline
C7 & If available Link to developer documentation/manual & Local docs under papers/opentelemetry-to-eeoap/. TODO: add public link after release. \\
\hline
C8 & Support email for questions & TODO: confirm public support contact before submission. \\
\hline
\end{tabular}
\caption{Code metadata (mandatory)}
\end{table}

\section{Motivation and significance}

Agent applications often produce traces, callback logs, span trees, or
tool-call records. These records are useful to developers, but they are not
automatically portable evidence objects. This package provides a small
executable bridge between OpenTelemetry-style telemetry and EEOAP-compatible
operation evidence. It does not define a new profile or move the validator
target. A reviewer can inspect a fixture, run the adapter, inspect the generated
statement and adapter report, and rerun the scoped tests.

\section{Software description}

The main adapter is tools/opentelemetry\_to\_eeoap\_adapter.py. Broader package
code and validator support live under agent\_evidence/. Inputs are stored under
examples/opentelemetry/. Generated EEOAP-compatible statements and adapter
reports are stored under generated/. Tests are in
tests/test\_opentelemetry\_to\_eeoap\_adapter.py.

The adapter parses trace spans, selects one accountable agent span, extracts
agent identity and operation metadata, checks parent-child span consistency,
resolves reachable tool execution spans, preserves source provenance, emits an
EEOAP-compatible statement, and uses the existing EEOAP validator path. Failure
diagnostics cover missing agent span, unresolved tool span, broken parent span
relation, and missing operation name. The adapter does not change the EEOAP
schema.

\section{Illustrative examples}

The valid-agent-trace fixture represents a research.answer operation with one
root agent span and two direct child tool spans. The valid-agent-workflow-trace
fixture represents workflow.execute with a deeper parent-child pattern. Both
generated statements pass the validator with ok=true and issue_count=0. The
version 1.16 clean-clone verification confirms the version 1.15 support package
checksum, scoped pytest, and validator results.

\section{Impact}

The adapter shows how trace-like telemetry can become a validator-ready
evidence object through a small, inspectable transformation package. It is
useful for artifact evaluation because reviewers can inspect source traces,
generated statements, reports, and tests without depending on a live telemetry
service. It also provides a baseline for later runtime integrations whose
claims must be tied to committed inputs, outputs, tests, and validation results.

\section{Conclusions}

The OpenTelemetry-to-EEOAP Adapter is a bounded research-software package for
mapping controlled OpenTelemetry-style agent trace fixtures into
EEOAP-compatible operation accountability statements. It is not a legal
accountability system, production integration, broad OpenTelemetry compatibility
claim, or new EEOAP profile. DOI, GitHub Release, pushed tags, final public
metadata, final references, and final SoftwareX submission remain TODO.

\section*{Limitations}

The fixtures are synthetic and only two valid trace contexts are included. The
package does not include real LangChain runtime integration, OpenTelemetry
Collector integration, broad OpenTelemetry implementation compatibility, legal
accountability proof, production readiness, or agent-output correctness.

\section*{Reproducibility and artifact availability}

The release-candidate branch is softwarex-otel-eeoap-release-candidate. The
version 1.15 support package supersedes the historical version 0.5 frozen
package as the forward SoftwareX support package. Version 1.16 verified this
package from a clean checkout. Local tags and metadata drafts exist, but they
are not public release metadata. DOI creation, GitHub Release creation, public
tag push, final release URL, and final artifact availability wording remain
TODO before formal submission.

\section*{Declaration of generative AI and AI-assisted technologies in the manuscript preparation process}

OpenAI ChatGPT/Codex was used for planning, drafting, code scaffolding,
documentation structuring, local command execution support, and review support.
The human author remains responsible for the content, code, claims, artifacts,
citations, limitations, and conclusions. The AI system is not an author. TODO:
adapt this wording to the target venue policy before submission.

\section*{Declaration of competing interest}

TODO: confirm final wording before submission. Draft statement: the author
declares no known competing financial interests or personal relationships that
could have appeared to influence the work reported in this paper.

\section*{Funding}

TODO: confirm final wording before submission. Draft statement: this research
received no specific grant from any funding agency in the public, commercial, or
not-for-profit sectors.

\section*{Data availability}

No human subject data, patient data, private production telemetry, or private
operational logs are used. Fixtures are synthetic OpenTelemetry-style trace JSON
files under examples/opentelemetry/. Generated statements and adapter reports
are under generated/. TODO: update after public release or archive.

\section*{References}

\begin{thebibliography}{00}
\bibitem{otel-genai} OpenTelemetry. Semantic conventions for generative AI systems. TODO: verify URL and status immediately before submission.
\bibitem{json-schema} JSON Schema Draft 2020-12. TODO: adapt format to SoftwareX final references.
\bibitem{prov-o} W3C PROV-O: The PROV Ontology. TODO: decide final provenance references.
\bibitem{eeoap-artifact} agent-evidence repository. TODO: replace with immutable release, archive, tag, or DOI before external submission.
\end{thebibliography}

\section*{Current executable software version}

TODO: complete after release metadata, tag, and public repository strategy are resolved.

\end{document}
